% ============================================================
%  第1章  绪  论
% ============================================================
\chapter{绪\ \ 论}

\section{研究背景与意义}
脑卒中（Stroke）是全球范围内导致成人残疾的主要原因之一\cite{卒中报告2023}。据统计，约$60\%$以上的脑卒中幸存者伴有不同程度的眼动功能障碍，包括扫视障碍、平滑追踪异常及空间忽略等\cite{eyemovement2021}。传统的眼动康复训练依赖于治疗师一对一指导或大型康复设备，存在成本高、场地受限、训练频次不足等问题\cite{saccade2020}。因此，开发一种便携、低成本、可居家使用的头戴式眼部康复装置具有重要的临床意义和应用价值。

\section{国内外研究现状}
在眼动康复领域，现有的康复训练系统主要分为两大类：一是基于大型眼动追踪设备的康复系统，如虚拟现实（VR）结合眼动追踪的训练平台\cite{pursuit2019}，这类系统精度高但体积大、成本昂贵；二是基于便携式设备的康复方案，如头戴式激光引导训练装置，利用激光投射光点引导患者进行注视、扫视和平滑追踪训练\cite{eyemovement2021}，但目前此类设备多以单片机为核心，缺乏姿态自适应能力和语音交互功能。

在姿态检测方面，基于微机电系统（MEMS）惯性测量单元（IMU）的头部姿态估计技术已成为康复监测领域的研究热点\cite{mahony2008}。其中，卡尔曼滤波（Kalman Filter）及其衍生算法被广泛用于多传感器数据融合\cite{kalman1960}，Madgwick等人提出的梯度下降算法也取得了良好的效果\cite{madgwick2010}。

针对康复训练中的视觉反馈问题，Bates等人的研究表明，平滑追踪训练能有效改善患者的注视稳定性\cite{pursuit2019}。在空间忽略训练方面，Rowe等人的系统综述指出，通过左右视野交替引导训练可以显著改善空间忽略症状\cite{saccade2020}。

\section{本文主要工作与章节安排}
本文设计并实现了一款基于STM32H7的脑卒中眼部康复头戴式装置。主要工作包括：\textcircled{1}设计了基于BMI088的头部姿态检测系统，实现了传感器坐标系到头部坐标系的轴系重映射，并融合了卡尔曼滤波与纯陀螺仪积分的姿态解算算法；\textcircled{2}提出了五种临床眼动训练模式的状态机架构，涵盖注视稳定性、扫视、平滑追踪、视觉聚焦和空间忽略训练；\textcircled{3}集成了语音交互模块，实现了训练过程的语音控制和音频反馈；\textcircled{4}设计了实时姿态安全监测机制，确保训练过程的安全性。

论文共分为六章：第一章绪论，介绍研究背景、现状和主要工作；第二章系统总体方案设计；第三章硬件设计与实现；第四章软件设计与实现；第五章系统测试与结果分析；第六章总结与展望。
